\section{Fibonacci residues: why the pattern must repeat}
\label{sec:fibonacci-periodicity}

Set $F_0=0$, $F_1=1$, and $F_{n+2}=F_{n+1}+F_n$. Modulo an integer
$m\ge2$, the next term depends on the previous \emph{two} terms.
The appropriate state is therefore
\[
 S_n=(F_n,F_{n+1})\pmod m,
 \qquad T(a,b)=(b,a+b)\pmod m.
\]
There are only $m^2$ pairs, so a state must eventually repeat.
But finiteness alone only gives eventual periodicity. Here we can do better:
\[
 T^{-1}(u,v)=(v-u,u)\pmod m.
\]
Every step has a unique predecessor, using only subtraction.

\begin{theorem}[Fibonacci periodicity from the starting pair]
For every $m\ge2$, the Fibonacci sequence modulo $m$ is periodic from
the start. Its least period $\pi(m)$ is the smallest positive $t$ for which
\[
 (F_t,F_{t+1})\congr(0,1)\pmod m.
\]
In particular, $\pi(m)\le m^2-1$.
\end{theorem}

\begin{proof}
If $S_i=S_j$ with $i<j$, apply the inverse step $i$ times to get
$S_0=S_{j-i}$. Repeating the forward step now gives a period beginning
at $S_0$, not just after an initial segment. Conversely, a period of the
sequence must return both $F_0$ and $F_1$, so it returns the starting pair.
Before the first return all states are distinct. The pair $(0,0)$ is fixed
by $T$ and has no other predecessor, so the orbit of $(0,1)$ never visits
it. There are at most $m^2-1$ available states.
\end{proof}

\begin{example}[A zero is not yet a restart]
Modulo $3$, the successive states are
\[
\begin{array}{c|rrrrrrrrr}
n&0&1&2&3&4&5&6&7&8\\\hline
F_n&0&1&1&2&0&2&2&1&0\\
F_{n+1}&1&1&2&0&2&2&1&0&1
\end{array}
\]
The pair first returns at $n=8$, so $\pi(3)=8$.
Although $F_4\congr0$, its successor is $2$, not $1$.
Looking only for a repeated Fibonacci value would give the wrong answer.
\end{example}

\begin{proposition}[CRT combines periods]
If $\gcd(m,n)=1$, then
\[
 \pi(mn)=\operatorname{lcm}(\pi(m),\pi(n)).
\]
\end{proposition}

\begin{proof}
CRT says that a pair returns modulo $mn$ exactly when it returns modulo
both $m$ and $n$. The return times modulo $m$ are precisely the multiples
of $\pi(m)$: divide any return time by $\pi(m)$ and use the minimality of
the remainder. Likewise for $n$. The first common return time is their
least common multiple.
\end{proof}

For example, the orbit modulo $2$ has length $3$, so
$\pi(6)=\operatorname{lcm}(3,8)=24$.
The connection is operational: CRT reassembles the two coordinates of
every state as well as the arithmetic on the clocks.

\begin{labbox}
For a small modulus, start at $(0,1)$, repeatedly replace $(a,b)$ by
$(b,a+b\bmod m)$, and stop on the first return to $(0,1)$.
The theorem proves this procedure terminates and that the answer is exact.
Record pairs, not just zeros. These small experiments need no extension fields.
\end{labbox}

\subsection*{What we can prove now, and what comes next}

The reversible-pair argument proves periodicity for \emph{every} modulus,
prime or composite. The estimate $\pi(m)\le m^2-1$ is an upper bound;
it does not assert that $\pi(m)$ divides $m^2-1$.
For example, direct iteration gives $\pi(5)=20$, and $20\ndivides24$.

A structural computation uses the roots of $X^2-X-1$.
Their $\sqrt5$ exists in $\F_p$ for $p=11$, but not for $p=7$.
Appendix~\ref{app:finite-fields} first builds arithmetic in an extension
field, then uses Frobenius and multiplicative orders to compute exact
periods. Those tools belong to a later unit; pair iteration and CRT
require only the arithmetic developed here.
